% !TEX root = ../main.tex
% File: chapters_part1/chap7_1.tex
% Nội dung cho Chương 7, Phần 1


\section{Các metric kinh điển}
\label{sec:classic_metrics}

Mỗi metric trong mục này đo một khía cạnh cụ thể của hiệu năng mô hình.

\subsection{Perplexity: đo chất lượng của mô hình ngôn ngữ}
\label{ssec:perplexity}
\textbf{Perplexity} (PPL), hay độ bối rối, đánh giá chất lượng nội tại của một \textbf{mô hình ngôn ngữ (Language Model)}, tức mô hình được huấn luyện để dự đoán từ tiếp theo. Nó định lượng mức “ngạc nhiên” (surprise) của mô hình khi gặp một câu văn thực tế: perplexity càng thấp, mô hình càng tốt.

Perplexity quan hệ trực tiếp với hàm mất mát Cross-Entropy. Với một chuỗi văn bản $W = (w_1, w_2, \dots, w_N)$:
\begin{equation}
    \text{PPL}(W) = \sqrt[N]{\frac{1}{P(w_1, w_2, \dots, w_N)}} = \left( \prod_{i=1}^{N} \frac{1}{P(w_i|w_1, \dots, w_{i-1})} \right)^{1/N}
    \label{eq:perplexity}
\end{equation}
Về mặt toán học, nó là nghịch đảo của xác suất trung bình trên mỗi từ, lấy căn bậc $N$.

Perplexity có hai hạn chế. Nó chỉ đo khả năng dự đoán xác suất, không trực tiếp đo chất lượng ngữ nghĩa hay sự mạch lạc của văn bản được sinh ra. Và nó rất nhạy với kích thước từ vựng, nên không thể so sánh perplexity của hai mô hình có từ vựng khác nhau.

\subsubsection{Perplexity chính là hàm mất mát, chỉ viết theo đơn vị khác}
Lấy logarit tự nhiên hai vế của công thức~\eqref{eq:perplexity}:
\[
\ln \text{PPL}(W) = -\frac{1}{N}\sum_{i=1}^{N} \ln P(w_i \mid w_1, \dots, w_{i-1}) = L ,
\]
vế phải chính là \textbf{cross-entropy trung bình trên mỗi token} -- đúng cái loss ta nhìn chạy trên biểu đồ khi huấn luyện. Vậy
\[
\boxed{\ \text{PPL} = e^{L}\ }
\]
Loss $2{,}0$ tương ứng PPL $\approx 7{,}39$; giảm loss được $0{,}1$ nghĩa là PPL được nhân với $e^{-0{,}1} \approx 0{,}905$, tức giảm khoảng $9{,}5\%$. Đây là đẳng thức hữu ích nhất cần nhớ về perplexity: \emph{đường loss và đường log-perplexity là một đường}.

\paragraph{Đọc perplexity như “số lựa chọn hiệu dụng”.} Một mô hình đoán đều trên từ vựng $V$ token có $P = 1/V$ ở mọi bước, nên PPL $= V$, đúng bằng kích thước từ vựng. Tổng quát, PPL $= k$ nghĩa là mô hình bối rối ngang với việc phải chọn đều giữa $k$ khả năng. Một LLM hiện đại đạt PPL cỡ một chữ số tới hơn chục trên văn bản web với từ vựng $128\,000$ token -- tức ở mỗi bước nó đã thu hẹp $128$ nghìn khả năng xuống còn tương đương chục khả năng.

\paragraph{Bit, và mối liên hệ với nén dữ liệu.} Dùng $\log_2$ thay cho $\ln$, ta được \textbf{bit trên mỗi token}, và PPL $= 2^{\text{bits}}$. Con số này có ý nghĩa vật lý rất cụ thể: theo lý thuyết mã hóa, đó chính là số bit trung bình cần để mã hóa một token nếu ta dùng mô hình làm bộ nén. Một mô hình ngôn ngữ tốt hơn \emph{đúng nghĩa đen} là một bộ nén tốt hơn -- đây không phải phép ẩn dụ mà là cùng một đại lượng nhìn từ hai phía.

\paragraph{Perplexity không so sánh được giữa hai tokenizer khác nhau.} Perplexity được tính \emph{trên mỗi token}, mà “token” là gì thì do tokenizer quyết định. Một tokenizer có từ vựng lớn cắt văn bản thành \emph{ít} token hơn, mỗi token mang nhiều thông tin hơn và do đó khó đoán hơn -- PPL cao hơn, dù mô hình chẳng hề tệ hơn. Ngược lại, tokenizer cắt vụn (ví dụ theo ký tự) cho PPL thấp một cách giả tạo. Muốn so sánh hợp lệ thì phải chuẩn hóa về một đơn vị không phụ thuộc tokenizer, phổ biến nhất là \textbf{bit trên mỗi byte} (bits-per-byte, BPB) \cite{gao2020pile}:
\[
\text{BPB} = \frac{N_{\text{token}}}{N_{\text{byte}}} \cdot \frac{L}{\ln 2} .
\]
Đây là lý do các báo cáo nghiêm túc về mô hình nền dùng BPB chứ không dùng PPL thô.

\paragraph{Hai cái bẫy còn lại khi báo cáo perplexity.}
\begin{itemize}
    \item \textbf{Phụ thuộc vào cách cắt cửa sổ ngữ cảnh.} Văn bản dài hơn cửa sổ của mô hình phải được cắt thành từng đoạn. Nếu cắt rời rạc, những token đầu mỗi đoạn bị đánh giá trong tình trạng gần như không có ngữ cảnh và PPL bị đội lên. Cách chuẩn là dùng \textbf{cửa sổ trượt} với bước nhảy nhỏ hơn độ dài cửa sổ và chỉ tính loss cho phần token mới. Cùng một mô hình, hai cách đánh giá này có thể cho những con số chênh nhau đáng kể -- nên luôn phải nói rõ đã đo thế nào.
    \item \textbf{PPL thấp không đồng nghĩa mô hình hữu ích hơn.} Sau khi căn chỉnh theo sở thích con người (mục~\ref{ssec:rlhf}), perplexity của mô hình trên văn bản web thường \emph{tăng lên} trong khi mức hữu ích với người dùng tăng rõ rệt. Perplexity đo mức khớp với phân phối dữ liệu, không đo mức hữu ích -- và sau RLHF thì hai thứ đó \emph{cố tình} không còn trùng nhau nữa.
\end{itemize}

\subsection{Precision, Recall và F1-score: nền tảng của các bài toán phân loại}
\label{ssec:f1_score}
Ba metric này được dùng cho phần lớn các bài toán \textbf{phân loại (classification)} và \textbf{trích xuất thông tin (information extraction)} như NER.

\subsubsection{Các thành phần cơ bản}
Khi phân loại một điểm dữ liệu có bốn kết quả:
\begin{itemize}
    \item \textbf{True Positive (TP):} Dự đoán là Dương tính, thực tế là Dương tính (dự đoán đúng).
    \item \textbf{True Negative (TN):} Dự đoán là Âm tính, thực tế là Âm tính (dự đoán đúng).
    \item \textbf{False Positive (FP):} Dự đoán là Dương tính, thực tế là Âm tính (lỗi loại I, “báo động giả”).
    \item \textbf{False Negative (FN):} Dự đoán là Âm tính, thực tế là Dương tính (lỗi loại II, “bỏ sót”).
\end{itemize}

\begin{center}
\bookimage[0.6\textwidth]{confusion_matrix}{Bảng hai hàng hai cột. Cột là nhãn dự đoán (Negative, Positive), hàng là nhãn thực tế (Negative, Positive). Bốn ô lần lượt chứa: số lượng TN (thực tế Âm, dự đoán Âm), FP (thực tế Âm, dự đoán Dương), FN (thực tế Dương, dự đoán Âm) và TP (thực tế Dương, dự đoán Dương).}
\captionof{figure}{Ma trận nhầm lẫn (Confusion Matrix) minh họa bốn kết quả của một bài toán phân loại nhị phân.}
\label{fig:confusion_matrix}
\end{center}

\subsubsection{Định nghĩa các metric}
\paragraph{Precision (độ chính xác)} Trong số tất cả những gì mô hình dự đoán là Dương tính, bao nhiêu phần thực sự là Dương tính:
$$ \text{Precision} = \frac{\text{TP}}{\text{TP} + \text{FP}} $$
Tối đa hóa Precision là giảm thiểu số báo động giả (FP). Đó là mục tiêu của các ứng dụng như phát hiện spam: thà bỏ sót một email spam còn hơn đưa một email quan trọng vào hòm thư rác.

\paragraph{Recall (độ phủ, hay độ nhạy)} Trong số tất cả những cái thực sự là Dương tính, mô hình tìm thấy được bao nhiêu:
$$ \text{Recall} = \frac{\text{TP}}{\text{TP} + \text{FN}} $$
Tối đa hóa Recall là giảm thiểu số bỏ sót (FN). Đó là mục tiêu của các ứng dụng như chẩn đoán y tế: thà chẩn đoán nhầm một người khỏe mạnh là có bệnh còn hơn bỏ sót một người thực sự có bệnh.

\paragraph{F1-score (điểm F1)} Precision và Recall thường đánh đổi nhau: một mô hình “cẩn trọng” có thể có Precision cao nhưng Recall thấp, và ngược lại. F1-score gộp cả hai thành một chỉ số duy nhất bằng trung bình điều hòa (harmonic mean):
$$ \text{F1-Score} = 2 \cdot \frac{\text{Precision} \cdot \text{Recall}}{\text{Precision} + \text{Recall}} $$
Trung bình điều hòa chỉ cao khi cả Precision và Recall đều cao, nên F1-score được dùng rộng rãi trong các tác vụ như NER, nơi ta quan tâm đến cả việc không nhận dạng sai và không bỏ sót thực thể.

Khi đánh giá trên nhiều lớp, người ta thường tính \texttt{macro-F1} (tính F1 cho mỗi lớp rồi lấy trung bình) hoặc \texttt{micro-F1} (tính tổng TP, FP, FN trên tất cả các lớp rồi mới tính F1).

\subsection{BLEU: đánh giá dịch máy}
\label{ssec:bleu}
\textbf{BLEU} (Bilingual Evaluation Understudy) được dùng chủ yếu cho \textbf{dịch máy (Machine Translation)}, đôi khi cho các tác vụ sinh văn bản khác. Nó dựa trên giả định rằng một bản dịch tốt có nhiều N-gram trùng lặp với các bản dịch tham khảo chất lượng cao do con người tạo ra.

Điểm BLEU được tính từ hai yếu tố. Thứ nhất là \textbf{độ chính xác N-gram đã sửa đổi (Modified N-gram Precision)}: tỷ lệ các N-gram (thường từ 1-gram đến 4-gram) trong câu do máy dịch tạo ra cũng xuất hiện trong các câu tham khảo; phần “sửa đổi” ngăn mô hình lặp lại một từ đúng nhiều lần để gian lận điểm. Thứ hai là \textbf{hình phạt cho câu ngắn (Brevity Penalty, BP)}: câu do máy dịch tạo ra mà ngắn hơn đáng kể so với các câu tham khảo sẽ bị phạt, nhằm ngăn mô hình chỉ sinh ra một vài từ rất chính xác để có Precision cao.
$$ \text{BLEU} = \text{BP} \cdot \exp\left(\sum_{n=1}^{N} w_n \log p_n\right) $$
Trong đó $p_n$ là độ chính xác của n-gram, và $w_n$ thường là $1/N$.

BLEU có ba hạn chế. Nó chỉ xét Precision, không xét Recall, nên không phạt một bản dịch bỏ sót thông tin. Nó không hiểu ngữ nghĩa, nên một bản dịch dùng từ đồng nghĩa hoàn hảo vẫn bị điểm thấp nếu không khớp chính xác N-gram. Và nó không đánh giá tốt các bản dịch có trật tự từ hay cấu trúc câu khác biệt nhưng vẫn đúng.

\subsection{ROUGE: đánh giá tóm tắt văn bản}
\label{ssec:rouge}
\textbf{ROUGE} (Recall-Oriented Understudy for Gisting Evaluation) được dùng chủ yếu cho \textbf{tóm tắt văn bản (Summarization)}. Nó đảo chiều so với BLEU: một bản tóm tắt tốt phải bao phủ được các N-gram quan trọng có trong các bản tóm tắt tham khảo của con người.

ROUGE là một họ metric, trong đó hai biến thể phổ biến nhất là:
\begin{itemize}
    \item \textbf{ROUGE-N:} Dựa trên sự trùng lặp N-gram, nhưng tính theo Recall.
        $$ \text{ROUGE-N} = \frac{\text{Số N-gram trong tham khảo cũng xuất hiện trong tóm tắt máy}}{\text{Tổng số N-gram trong tham khảo}} $$
        Thường dùng \texttt{ROUGE-1} (unigram) và \texttt{ROUGE-2} (bigram).
    \item \textbf{ROUGE-L:} Dựa trên chuỗi con chung dài nhất (Longest Common Subsequence, LCS). Nó không yêu cầu các từ phải liền kề nhau, nên nắm bắt được sự tương đồng về trật tự từ tốt hơn.
\end{itemize}
Giống BLEU, ROUGE chỉ dựa trên sự khớp bề mặt của từ và không xét ngữ nghĩa: một bản tóm tắt diễn đạt lại ý tưởng hoàn toàn chính xác vẫn có thể bị điểm ROUGE thấp.

Dù có những hạn chế trên, các metric kinh điển vẫn được dùng rộng rãi vì chúng tự động, nhanh và cho một thước đo khách quan (dù không hoàn hảo) để so sánh các mô hình trong quá trình phát triển.